\documentclass{DISERTASI-ITS}
% tambahkan package yang dibutuhkan disini

\numberwithin{equation}{chapter}
\theoremstyle{plain}
\newtheorem{thm}{Teorema}[section]
\newtheorem{prop}[thm]{Proposisi}
\newtheorem{lemma}[thm]{Lemma}
\newtheorem*{cor}{Akibat}
\theoremstyle{remark}
\newtheorem{rmk}{Remark}
%\newtheorem*{prf}{Bukti}
\renewcommand{\qedsymbol}{$ \blacksquare $}
\theoremstyle{definition}
\newtheorem{defi}{Definisi}[section]
	
\usepackage{datetime}
	\newdateformat{monthdayyeardate}{%
	\THEDAY~\monthname[\THEMONTH]~\THEYEAR}

\usepackage{caption} 
\captionsetup{margin=0pt,singlelinecheck=off,format=plain,indention=.15cm,justification=raggedright}

\begin{document}
        \Keterangan{proposal} % kalau kosong berarti Keterangan pada lembar Judul diberi keterangan Disertasi, selain itu berarti Proposal Disertasi, 
	\Judul{Judul Bahasa Indonesia} 	% Tuliskan Judul Yang sesuai (bhs.Indonesia)
	\JudulEng{Judul Bahasa Inggris} 				% Tuliskan Judul Yang seuai (bhs.Inggris)
	\Nama{Nama Mahasiswa} 							% Tuliskan Nama penulis proposal tesis
	\NRP{XXXXXXXX} 							% Tuliskan Nrp penulis proposal tesis
	\Jurusan{Matematika}						% Tuliskan Departement Penulis tesis
	\Department{Mathematics}					% Tuliskan Departement Penulis tesis (bhs. Inggris)
	\BidangStudi{Analisis dan Aljabar} 			% Tuliskan bidang sesui minat penulis
	\Tahun{2022} 								% Tuliskan tahun pembuatan proposal tesis
	\Fakultas{Sains dan Analitika Data} 		% Tuliskan Fakultas pembuatan proposal tesis
	\Faculty{Science and Analytical Data} 		% Tuliskan Fakultas pembuatan proposal tesis (bhs. Inggris)
	
	\KepalaDep{Subchan, S.Si., M.Sc., Ph.D.}
	\NipKepalaDep{197105131997021001}
	\Pembimbing{Prof. Dr. Subiono, M.S.} 		% Tuliskan nama pembimbing 1
	{Dr. Dieky Adzkiya, M.Si.} 					% (Dr. Dieky Adzkiya, M.Si.) Tuliskan nama pembimbing2, bila hanya satu jangan diisi
	{}          								% Tuliskan nama pembimbing3, bila hanya dua jangan diisi
	\NIPPembimbing{195704111984031001} 			% Isi dengan NIP pembimbing 1
	{198305172008121003} 						% (198305172008121003) Isi dengan NIP pembimbing kedua bila ada
	{}    										% Isi dengan NIP pembimbing ketiga bila ada
	
	\Penguji{Dr. Mahmud Yunus, M.Si.}   	% Isi nama penguji  1  yang sesuai
	{Dr. Dra. Rinurwati, M.Si.}           	% Isi nama penguji  2  yang sesuai
	{Dr. Dwi Ratna S, S.Si., M.T.}           	% Isi nama penguji  3  yang bila ada bila tidak kosongkan
	%	
		\NIPPenguji{196204071987031005}  	% Isi Nip penguji 1 yang sesuai
		{196403041989032002}             	% Isi Nip penguji 2 yang sesuai
		{196904051994032003}             	% Isi Nip penguji 3 yang sesuai bila ada bila tidak kosongkan
	%	
	
	\TanggalDisetujui{6 Desember 2020} %  \monthdayyeardate\today   %:tanggal pelaksanaan seminar proposal (21 September 2020)  5 Desember 2020
    \TanggalSeminar{4 Januari 2023}
    \HariSeminar{Rabu}
    \TempatSeminar{Ruang Sidang II Departemen Matematika, Gedung U Lantai 2}
    \TanggalUjian{11 Februari 2023}
    \PeriodeWisuda{Dua}
	
	\BagianAwal
	
	\LembarJudul  % menampilkan halaman judul
	\TitlePage
	%\LembarPersetujuan % Menampilkan Lembar Persetujuan Seminar oleh Promotor
	\LembarPengesahan % Menampilkan Lembar Pengesahan Seminar Proposal oleh promotor dan penguji
	%\LembarPengesahanDisertasi % Menampilkan Lembar Pengesahan Disertasi
	\begin{Abstrak} % Abstrak bhs. Indonesia
		\input{AbstrakInd/AbstrakInd} % isikan foder dan file abstrakind.tex
	\end{Abstrak}
	
	\begin{Abstract} % Abstrak bhs. Inggris
		\input{AbstrakEnglish/AbstrakEnglish} %  isikan folder dan file abstrakenglish.tex
	\end{Abstract}
	\DaftarIsi % menampilkan secara otomatis semua Daftar Isi tesis
	\DaftarGambar
	\DaftarTabel
	\BagianInti
	\input{BabI/BabI} % isikan nama folder dan file BabI.tex sebagai Bab pendahuluan

	\input{BabII/BabII} % isikan nama folder dan file BabII.tex  
	\input{BabIII/Bab3} % isikan nama folder dan nama file Bab3.tex
	
	\DaftarPustaka{ReferensiDisertasi}	
	
\end{document}
